\subsection*{Grading and the attractive sign}
Choose finite positive-mass, positive-clock, positive-ruler polar packets
$A=UE_0$, $E_1=UE_0U^\dagger$ and
$B_j=E_j^{1/2}(N\psi)E_j^{1/2}>0$.
There are 2 canonical realizations:
\[
 h_0=\begin{pmatrix}0&B_0U^\dagger\\UB_0&0\end{pmatrix},
 \qquad
 h_1=\begin{pmatrix}0&U^\dagger B_1\\B_1U&0\end{pmatrix}.
\]
Their filled-band energies are $-\operatorname{Tr}B_0$ and
$-\operatorname{Tr}B_1$.
Putting the compensating boson oscillators in the opposite channel
gives respectively
$\pm\operatorname{Tr}[(N\psi)(E_1-E_0)]$.
Both realizations have positive norms and nonnegative excitation
forms: their one-body eigenvectors are $(f,\pm Uf)$, and independently
$h_0^2=\operatorname{diag}(B_0^2,UB_0^2U^\dagger)$, with the
corresponding reversed formula for $h_1$.
A number-conserving Hermitian dressing which kills excitation
numbers 0 and 1 keeps $F^\dagger BF\ge0$ in either realization.
[X: band isometries; squared-block spectral calculus.]

Reverse every packet orientation in the mass rectangle.
This changes the law's nonuniform couplings but preserves its
uniform readout, the 2 mass sum rules, positive source forms and
constant-allocation null identity.
It sends $W_0\mapsto-W_0$ and $k\mapsto-k$, where
$W_0^{[2]}=-k\int\nabla\psi\cdot\nabla(N\psi)$.
The canonical channel assignments above are an explicit
counterfamily to selection by positive norm or spectral boundedness.
They are not a mere sign relabelling of an unchanged Hamiltonian.

The constrained boundary bound also admits both orientations.
On a finite open graph let $w_x$ be the signed vacuum density,
$\sum_xw_x=0$, and impose the interior clock rows
$C_x=\psi_xw_x+e_x=0$, with $e_x\ge0$.
At a unit collar, retaining its source, the same-law charge is
\[
 E_{\rm bdy}=\sum_{x\in{\rm collar}}(w_x+e_x)
            =\sum_{\rm all\ x}\frac{e_x}{\psi_x}\ge0 .
\]
Alternatively the comparison lapse $N_x=\psi_x^{-1}$ makes the
vacuum trace vanish and gives the same source sum.
[X: weighted collar sum; constant-allocation comparison lapse.]
This bound assumes an actual solution and the stated open exterior.
It proves neither existence nor an unrestricted off-constraint
lower bound; varying a positive lapse against a negative vacuum
clock coefficient can make the unreduced energy unbounded.

A concrete additional selecting principle is attraction of a weak
positive rest source in the declared clock chart.
At continuum leading order its rows give
$k\Delta f+\rho=0$ and $n=-2f$, with
$\psi=1+f$, $N=1+n$ and vanishing perturbations at infinity.
Thus
\[
 n(x)=-\frac1{2\pi k}\int
       \frac{\rho(y)}{|x-y|}\,d^3y .
\]
The test energy $m(1+n)$ attracts toward the source precisely for
$k>0$; the reversed law gives the opposite force.
$G$ is defined by Cavendish measurement.
The selected inventory must satisfy its single Newton normalization;
energy positivity does not supply the attractive orientation.

\paragraph*{SOURCES.}
\begingroup\footnotesize
\path{<internal>/DERIVATION.md},
sections 1--2: explicit canonical channel assignments, Gram sources
and signed scalar weight.  The sign counterfamily is reconstructed
above without adopting a numerical inventory.
\path{<internal>/rederivation_final3_vacuum_surface.tex},
section 1, and
\path{<internal>/redteam_final3_vacuum_surface.tex},
boundary identity: finite-collar and comparison-lapse methods.
\path{<internal>/STATE.tex},
model definition and the source/bridge entry tagged [S] for the
unrestricted positive-clock energy; and
\path{<internal>/DECISIONS.md},
section (d): grading and attractive sign.
The unrestricted-energy statement and the constrained bound have
different domains; neither is an internal-ruler dynamical stability
theorem.
\endgroup
% Governing source: <internal>/GOAL_alt.md,
% duty 4(2). Counterfamily uses ordinary positive CCR/CAR realizations.
% The positive Gram source is e_x=psi_x sum_r Q_rx^dagger Q_rx
% in the prescribed commuting representative, with all collar terms.
% Weak-rest-source derivation: the vacuum ruler row is
% k Delta n+2k Delta f=0 and the leading rest-mass ruler derivative
% vanishes; this gives n=-2f under the isolated boundary conditions.
% Delta(1/|x-y|)=-4pi delta(x-y) fixes the displayed sign and factor.
